La réalisation du \textbf{Personal Behavior Anomaly Detection System (PBADS)} a constitué une étape majeure de notre parcours académique.

Ce projet nous a confrontés à un défi essentiel : concevoir une solution numérique capable d'analyser les comportements quotidiens des utilisateurs, de détecter automatiquement les anomalies à l'aide d'algorithmes d'apprentissage automatique, et de générer des recommandations personnalisées pour améliorer la qualité de vie.

Nous avons conçu, développé et mis en service un système fonctionnel basé sur une architecture microservices qui automatise la détection d'anomalies comportementales, génère des alertes intelligentes avec niveaux de sévérité, et fournit des recommandations personnalisées via l'intégration avec Google Gemini AI. L'application permet désormais une analyse en temps réel des habitudes de vie, un suivi précis des tendances comportementales, ainsi qu'une visualisation interactive des données via un tableau de bord moderne avec filtres temporels.

Tout au long de ce projet, nous avons su relever plusieurs défis techniques et organisationnels, notamment la modélisation des besoins utilisateurs avec les diagrammes UML, la conception d'une architecture microservices robuste avec Spring Boot, l'intégration de mécanismes d'authentification JWT et de contrôle d'accès basé sur les rôles, l'implémentation de la communication asynchrone avec Kafka et synchrone avec HTTP/REST, l'intégration d'algorithmes de machine learning (Isolation Forest, One-Class SVM) pour la détection d'anomalies, l'intégration avec Google Gemini AI pour la génération de recommandations, et l'optimisation de l'ergonomie des interfaces avec un thème sombre moderne et des graphiques interactifs.

Le respect des contraintes de délai, la rigueur dans la planification agile par itérations, et l'esprit collaboratif au sein de notre équipe ont été des facteurs clés dans l'atteinte de nos objectifs.

Cette expérience nous a permis de renforcer nos compétences en développement fullstack, en architecture microservices, en machine learning, ainsi qu'en intégration d'APIs d'intelligence artificielle.

Elle nous a aussi sensibilisés à l'importance d'une documentation claire, de cycles de test rigoureux, de la modélisation UML pour la compréhension des besoins, et de la prise en compte continue des retours utilisateurs pour garantir la qualité et l'évolutivité du produit.

À l'avenir, nous prévoyons d'enrichir le système avec des fonctionnalités supplémentaires telles qu'une analyse prédictive plus avancée, une intégration avec des dispositifs IoT (montres intelligentes, capteurs), une personnalisation accrue des recommandations basée sur l'historique utilisateur, ou encore une application mobile native pour un suivi encore plus pratique.

En résumé, cette réalisation a été une aventure formatrice et stimulante.

Nous sommes fiers du système que nous avons mis en place et restons motivés à poursuivre son amélioration pour répondre toujours mieux aux besoins d'une détection d'anomalies comportementales moderne, intelligente et accessible.
